\section{Background and Related Work}
\label{sec:background}
Remote shell protocols expose different notions of responsiveness to an \sout{AI-like}\textcolor{blue}{agent-like} terminal client. A session can be considered \sout{good}\textcolor{blue}{effective} because it becomes usable quickly, \textcolor{blue}{because short command responses arrive with low latency,} because typed characters become visible despite packet loss, or because command output is delivered as a complete byte stream. These goals are not equivalent. A protocol that minimizes visible delay may discard intermediate terminal states, while a protocol that preserves every byte may delay feedback until missing data has been retransmitted. This section summarizes the protocol mechanisms behind these trade-offs and relates them to prior measurement studies.

SSHv2 \cite{Ylonen2006RFC4251} is the de facto mechanism for secure remote terminal access. It runs over TCP and provides ordered, reliable delivery for encrypted SSH channels \cite{Bellare2002}. This byte-stream semantics is useful when the client must receive complete command output \textcolor{blue}{for programmatic parsing}, but it also inherits TCP \sout{HoL}\textcolor{blue}{head-of-line (HoL)} blocking: a lost segment can delay later data until retransmission completes \cite{Scharf2007}. For a closed-loop terminal client, this can appear as delayed command completion\textcolor{blue}{, delayed output availability,} or delayed visible feedback. SSHv2 setup also includes a TCP handshake followed by SSH key exchange and authentication, so the setup cost grows with \sout{Round Trip Time}\textcolor{blue}{round-trip time} (RTT) \cite{Ylonen2006RFC4253}. \textcolor{blue}{These properties make SSHv2 a natural reference point for byte-complete delivery, especially in workloads that stream large outputs through a single shell channel.}

SSH3 redesigns SSH over HTTP/3 and QUIC \cite{michel2023}. Its main architectural promise is lower setup latency: QUIC integrates transport and cryptographic negotiation, enabling 1-RTT setup and 0-RTT resumption for eligible connections \cite{Iyengar2021}. QUIC also removes HoL blocking across independent streams. For terminal workloads, however, the benefit depends on how shell traffic is mapped onto streams. If a workload consists of independent channels, QUIC stream multiplexing can prevent loss in one stream from blocking unrelated streams. If a workload produces one ordered terminal output flow, such as a large command printed through a single interactive shell, the stream still has to be completed in order before byte-complete output is available to the client. \textcolor{blue}{Consequently, SSH3 should be evaluated separately for session setup, short command-response loops, and large single-stream outputs rather than treated as uniformly faster than SSHv2.}

Mosh makes the opposite trade-off \cite{Winstein2012}. It uses SSH only to authenticate and bootstrap the session, then switches to a UDP channel carrying terminal state updates through the State Synchronization Protocol (SSP). SSP treats the terminal as a replicated \textcolor{blue}{screen} state and sends compact updates for the latest \sout{screen}\textcolor{blue}{visible terminal contents}. This design avoids waiting for every intermediate byte and makes lost datagrams superseded by newer state. Mosh further improves perceived responsiveness through predictive local echo. The cost is completeness: output that scrolls past between screen updates may never be delivered to the client. This is acceptable for many human editing sessions\textcolor{blue}{, where the latest visible screen is often sufficient}, but it is risky for automated clients that parse command output as data. \textcolor{blue}{Therefore, Mosh latency measurements must be interpreted together with output completeness; a fast observable completion is not necessarily a byte-complete completion.}

Prior work provides valuable but fragmented evidence. Mosh was originally evaluated for human-perceived responsiveness under mobile connectivity, showing that state synchronization and local echo can substantially reduce visible keystroke delay compared with SSHv2 \cite{Winstein2012}. QUIC deployment and measurement studies suggest similar expectations for SSH3: QUIC can improve selected latency and user-experience metrics at Internet scale \cite{Langley2017}, but later measurements show that these advantages depend on implementation choices, workload type, and network path \cite{Kakhki2017}. \textcolor{blue}{Recent QUIC throughput measurements further show that implementation behavior, stream concurrency, and user-space packet processing overhead can affect bulk transfer performance \cite{Konig2025}. These findings are important for remote shells because a terminal session can alternate between short control messages, large sequential outputs, and screen-redraw traffic.}

At the application layer, \textcolor{blue}{recent agent benchmarks show why these transport-level behaviors matter. }SWE-agent \sout{shows}\textcolor{blue}{demonstrates} that software-engineering agents \sout{often }interact with systems through an agent-computer interface that includes shell commands, file operations, and terminal observations \cite{Yang2024}. Terminal-Bench further illustrates that command-line environments are an important substrate for evaluating terminal-centered automation tasks \cite{merrill2026}. \textcolor{blue}{Other agent benchmarks similarly focus on whether an agent completes a task, not on whether the remote-shell transport delayed, reordered, or omitted the observations used by the agent \cite{Liu2024}.} These studies motivate the \sout{AI-like terminal interaction model considered in this paper}\textcolor{blue}{agent-like terminal workloads}, but they \sout{evaluate agent/task-level behavior rather than isolating the remote-shell transport layer}\textcolor{blue}{generally treat the shell connection as a reliable local interface rather than a networked feedback channel}.

As a result, the practical question is no longer which remote-shell protocol is universally fastest. Instead, it is which protocol best serves each terminal feedback pattern: session setup, small command responses, large byte-complete output, or visible editing feedback. What remains missing is a controlled comparison of SSHv2, SSH3, and Mosh on the same terminal-event workloads with \sout{both }latency and output completeness measured \textcolor{blue}{from the client side} under \sout{the same network scenarios}\textcolor{blue}{both controlled impairments and real network paths}. This paper fills that gap \sout{for AI-like interactive terminal clients}\textcolor{blue}{using simulated terminal workloads that isolate transport behavior from model reasoning and task-specific policy decisions}.
